\begin{center}
\LARGE
Thermodyanmics with Maple


\end{center}

\begin{center}
\Large
Ross Taylor
\end{center}

\noindent
Date:  July 18th (Thursday)\\
Time:  16:50-17:15\\

\begin{center}
\large
Abstract
\end{center}

Computer algebra systems (of which Maple is but one example) have enormous potential, not just
in thermodynamics, but in all areas of science and engineering. However, to be useful in
thermodynamics a computer algebra system needs: 

\begin{enumerate}

\item
The ability to work with total differentials of undefined functions, for example dS where
    $S=S(U,V)$.
\item
To be able to work with the subscripted (indexed) partial derivatives of thermodynamic
    functions found in many thermodynamic formulae.
\item
The ability to differentiate undefined or arbitrary sums (a summation where the upper limit
    is a symbol rather than a number). 

\end{enumerate}
As far as the author is aware, none of the commercially available computer algebra is able to do
any of the above right out of the box. In this paper we describe extensions to Maple so that it can
be used to rapidly develop the expressions needed to compute thermodynamic properties of
mixtures of any number of components. This paper also presents examples of thermodynamic
computations carried out using Maple with an emphasis on graphical visualization of both
symbolic and numerical results. One of the advantages of using a computer algebra system is the
greatly increased accuracy of derivations of thermodynamic properties and numerical
computations. That is, students are much more likely to get it right. On the debit side is a possible
decrease in the students ability to carry out such derivations and computations by hand. Other
advantages and disadvantages of using a computer algebra package in undergraduate education will
also be discussed. 

